\subsection{Aligning the two relative degree-three periods}

\begin{proposition}[A-L13: Proved; an exact simultaneous boundary criterion]
\label{AL13}
Keep the actual maps $f:X\to A$ and $g:X\to B$, the weight
$w$, and every hypothesis of AT15 with $q=3$.
Fix its already constructed data
\[
 \begin{gathered}
 z:D(2)\longrightarrow Z^3(X^w),\qquad
 a:K(1)\longrightarrow X^{w,2},\\
 b:D(2)\longrightarrow X^{w,2},\qquad
 da=z|_{K(1)},\qquad db=mz,\qquad b|_{K(1)}=ma.
 \end{gathered}
 \tag{AL.71}
\]
where $m\ne0$.
Write $\epsilon_f,\epsilon_g\in S_R^w/mS_R^w$ for the
actual residues of \textup{(AT.65)} for these fixed data.

For each closed relative one-form
$c:K(1)\to Z^1(X^w)$, let $Y(c)$ be the relative strict
PD algebra on $u(D(2))$ in form degree two and $v(D(2))$
in form degree one, with
\[
 du=z,\qquad dv=mu-b,\qquad
 u|_{K(1)}=a,\qquad v|_{K(1)}=c.
 \tag{AL.72}
\]
The two maps $f,g$ extend simultaneously to $Y(c)$ exactly when
\[
 \epsilon_f=\epsilon_g=\tau_X(c)\pmod {mS_R^w}.
 \tag{AL.73}
\]
Consequently the set of such simultaneous extensions, allowing
$c$ to vary, is nonempty exactly when $\epsilon_f=\epsilon_g$.
The proof constructs the boundary $c$, the algebra, and both
strict maps whenever this equality holds.

Each algebra $Y(c)$ is a pointwise free relative extension on
the two copies of $K(2)$, preserves simplices zero and one,
preserves $H^0$ and $H^1$, and is a rational quasi-isomorphism
in every simplex and weight. Its terms are torsionfree, and
remain free over $R$ when the terms of $X$ are free.
These are assertions about this actual relative repair; no
equality of the two residues for arbitrary input maps is asserted.
\end{proposition}

\begin{proof}
The defining PD ideal in \textup{(AL.72)} is preserved by the
differential. On the first boundary relation its differential
is $z|_{K(1)}-da=0$; on the second it is
$ma-b|_{K(1)}-dc=0$.
All maps and relations are homogeneous and simplicial.
The split inclusion $K(1)_n\subset D(2)_n$ with free quotient
$K(2)_n$ gives, by the free PD universal property, the claimed
pointwise free presentation and the injection $X\hookrightarrow
Y(c)$. In degrees zero and one the quotient parameter vanishes,
so the corresponding simplicial algebras equal those of $X$.

We give an actual isomorphism that reduces every boundary choice
to the previously proved repair. The element represented by $c$
is in $Z_1N(X^{w,1})$. The horizontal row bound from AT15
therefore supplies $\ell_2\in N_2X^{w,1}$ with $d_0\ell_2=c$.
Under Dold--Kan this is a simplicial linear map
\[
       \ell:D(2)\longrightarrow X^{w,1},
       \qquad \ell|_{K(1)}=c.
\]
Put $b'=b+d\ell$. It satisfies $db'=mz$ and
$b'|_{K(1)}=ma$, since $dc=0$.
Denote by $Y(0;b')$ the zero-boundary relative algebra of AT15
with $b'$ in place of $b$, and use $u_0,v_0$ for its generators.
There are mutually inverse strict simplicial PD-DGA maps over $X$
\[
 \begin{aligned}
 \Phi:Y(c)&\longrightarrow Y(0;b'),&
       u&\longmapsto u_0,&v&\longmapsto v_0+\ell,\\
 \Psi:Y(0;b')&\longrightarrow Y(c),&
       u_0&\longmapsto u,&v_0&\longmapsto v-\ell.
 \end{aligned}
 \tag{AL.74}
\]
They have the specified boundary values, and
\[
 d(v_0+\ell)=mu_0-b'+d\ell=mu_0-b,
 \qquad d(v-\ell)=mu-b-d\ell=mu-b'.
\]
The equations for $u,u_0$ also agree. The graded free PD
universal property supplies the maps on all products and divided
powers. On the generators their two composites are the identity,
which proves the stated isomorphism without identifying either
presentation by an unstated change of convention.

For completeness, the pointwise rational contraction can be read
directly from this isomorphism. In $Y(0;b')_{\mathbb Q}$ replace
$u_0$ by $u_0-b'/m$. Its bottom is zero and its differential is
zero, while $dv_0=m(u_0-b'/m)$ and $v_0$ has zero bottom.
This is the rational identity disk on each free quotient
parameter. The explicit disk contraction used in AT15 proves
the rational quasi-isomorphism. No new form-zero variable occurs,
so $H^0$ is unchanged. In form degree one the coefficient of
each independent new $u_0$ in the differential of a cycle is
$m$ times the corresponding coefficient of $v_0$.
Torsionfreeness makes every such coefficient zero. Thus all
one-form cycles and all their form-zero boundaries are old,
proving preservation of $H^1$ as well.

It remains to compute the exact effect on the two target
obstructions. For $f$, choose the actual maps $U,V$ in AT15,
so that
\[
 dU=fz,\qquad U|_{K(1)}=fa,\qquad dV=mU-fb.
\]
For the data $b'$ the maps $U$ and $V-f\ell$ satisfy the
same equations with $b'$ in place of $b$. Naturality of the
coefficient period, under the given coefficient-preserving map,
therefore gives the identity
\[
 \epsilon_f(b')
    =\tau_A\bigl(V|_{K(1)}-fc\bigr)
    =\epsilon_f(b)-\tau_X(c)\pmod {mS_R^w}.
 \tag{AL.75}
\]
The identical computation for $g$ gives
$\epsilon_g(b')=\epsilon_g(b)-\tau_X(c)$.
AT15 supplies actual maps out of $Y(0;b')$ precisely when
the respective residues vanish. Composing those maps with
$\Phi$ proves the sufficiency in \textup{(AL.73)}.
Conversely a map from $Y(c)$ sends $v$ to a primitive whose
bottom is $fc$ or $gc$. Its defining period is exactly the
corresponding value in \textup{(AL.73)}, proving necessity.

Finally suppose the two residues agree, and select a coefficient
$r\in S_R^w$ representing their common class.
Horizontal exactness in form degree zero supplies
$h\in N_1X^{w,0}$ with $d_0h=r$.
The closed one-form $c=dh$ has both faces zero, and
\[
                    \tau_X(c)=d_0h-d_1h=r.
 \tag{AL.76}
\]
It is therefore an actual admissible boundary choice for
\textup{(AL.73)}. Choose $\ell$, form $b'$, and use AT15
for the two zero residues, then compose with \textup{(AL.74)}.
This constructs both target maps to the same source algebra.
If any simultaneous choice existed while the original residues
were unequal, the necessity already proved would make both
equal to the same class $\tau_X(c)$, a contradiction.
\end{proof}
